\documentclass{clean_cv}

% Add a BibTeX-style file encoding all of your publications to include here. You can export this from Zotero. Only include
% of the publications you want to appear here!
\addbibresource{publications.bib}

\author{Luigi Pizza}
\headlineposition{College Student}

\begin{document}

\maketitle
% In this section, you can use any of the FontAwesome icons. The commands \faCenter and \faCenterStyle have been defined to properly center the icons
% when using the default font settings.
%
% You can use any of the icons listed in the fontawesome5 package documentation (https://ctan.math.utah.edu/ctan/tex-archive/fonts/fontawesome5/doc/fontawesome5.pdf)
% If you need to specify a specific style (as is done here for the address card), you should use the two-argument \faCenterCycle command
\vspace*{-5px}
\begin{center}
\begin{tabular}{llll}
\href{tel:+393668341405}{\faCenter{phone-alt} +39 3668341405} &
\href{mailto:luigi.pizza2002@gmail.com}{\faCenter{envelope} luigipizza.jobs@gmail.com} & 
\href{https://www.linkedin.com/in/luigipizza/}{\faCenter{linkedin} Luigi Pizza} &
\href{https://github.com/luigi-pizza}{\faCenter{github} luigi-pizza} 
\\
\end{tabular}
\end{center}

\vspace{-1.5em}

\section{Education}

% The datetabular environment takes one argument, which is the width of the left date column. As seen here:
%   9em is a good choice for "dual-date" formats (e.g. Sep 2015 - Nov 2019).
%   4em is a good choice for month/year dates (Sep 2014).
%   2em is a good choice for year-only dates (as seen in the publications)
\begin{datetabular}{4.2em}
% This is just a tabular environment, for the most part. The dateentry command has been defined for
% convienence. It takes two arguments, the first is the date and the second is whatever you wish placed to the right.

\dateentry{Sep 2021 Jul 2024}{
\textbf{Sapienza - Università di Roma} 

\textit{BS in \href{https://corsidilaurea.uniroma1.it/en/corso/2022/30786/home}{Computer Science and Artificial Intelligence}, GPA: 4.0 (29.89/30 Weighted Average)
Final Grade: 110/100 cum laude}}
% Relevant Courses:

% - Programming 1 (Python) \quad - Programming 2 (Java OOP) \quad - Computer Architecture

% - Algorithms and Data Structures \quad - Data Management and Analysis \quad - Systems and Networking

% - Artificial Intelligence and Machine Learning
% }


\dateentry{Sep 2021 Sep 2024}{
\textbf{\href{https://www.collegiocavalieri.it/}{Collegio Universitario dei Cavalieri del Lavoro "Lamaro Pozzani"}} 

The Collegio Universitario dei Cavalieri del Lavoro "Lamaro Pozzani" is a merit-based residence hall recognized by the "Ministero dell'Università e della Ricerca". Please refer to the \href{https://www.linkedin.com/in/luigipizza/}{LinkedIn page} for more information.}



\dateentry{Mar 2023 May 2023}{
\textbf{\href{https://cyberchallenge.it/}{CyberChallenge.IT}}

\textit{"the main Italian initiative to identify, attract, recruit, and place the next generation of CyberSecurity professionals."}
}

\end{datetabular}

\section{Extracurriculars}

\begin{datetabular}{2em}
\dateentry{2024}{
\textbf{\href{https://swerc.eu/2023/about/}{Southwestern Europe Regional Contest (SWERC) 2023-2024}}

I've represented "Sapienza - University of Rome" (team \href{https://swerc.eu/2023/theme/scoreboard/domjudge.lip6.fr/public/team/48.html}{SapienzaRed}) at the \href{https://swerc.eu/2023/theme/scoreboard/domjudge.lip6.fr/public.html}{SWERC competition}: a programming contest for teams of three students, focused on algorithmic problem solving and practical coding. 
}



\dateentry{2023}{
\textbf{\href{https://www.uniroma1.it/en/pagina/honours-programmes}{Honours Program}}

The Honours Program is integrated with the regular study plan and is addressed to students willing to deepen their knowledge and capabilities through extracurricular activities in selected topics.

I have been selected to take part in the honors program offered by the Computer Science Department. My selected Area of Interest is "Blockchain technology" under the supervision of Professor Alessandro Mei.
}

\dateentry{2024 2022}{
\textcolor{gray}{Participation to} \textbf{\href{https://www.euca.eu/}{EucA} Projects} \quad \textit{\href{https://www.linkedin.com/posts/luigipizza_eye2023-strasbourg-activecitizenship-activity-7140668305626902528-aroh}{EYE2023 Quiz Tournament}, \href{https://www.euchangemakers.com/genz-dubrovnik-croatia}{GenZ Dubrovnik}, \href{https://www.euchangemakers.com/egdszeged}{EGD Szeged}, \href{https://www.euchangemakers.com/egdrome}{EGD Rome}}

Projects developed by \href{https://www.euca.eu/}{EucA} and supported by the European Parliament, tackling civic participation, guerrilla marketing and cyber security.
}

\dateentry{2023}{
\textbf{\href{https://fa23.bici.events/}{FA'23:}} 
\textit{il Futuro Annunciato}

A brief workshop aiming to provide insight into the new frontiers of computer science through seminars, interactive panel discussions, and informal conversations with computer researchers from academia, industry leaders, Ph.D.'s, and startup founders.
}


\dateentry{2022}{
\textbf{\href{https://buca22.bici.events/home}{BUCA’22:}} \textit{Challenges in building Billion Users Cloud Applications}

Summer School revolving around the challenges of building billion-user systems and ways to address these challenges. The Instructors are 3 High-Ranking Google Engineers specialised in this field. 
}

\dateentry{2021}{
\href{http://olimpiadi.dm.unibo.it/2021/05/28/finale-nazionale-della-gara-individuale-tabellone-2021/}{Bronze Medal in the Italian Maths Olympiads}
}
% There is something seriously funky happening between the tabular commands and the end of the itemize blocks.
% The command \eatvspace should be used if extra space appears at the end of a datetabular environment.


\end{datetabular}

\section{Languages}

\begin{center}\\[7px]
\begin{tabular}{lllll}
\textbf{Italian}  & Native Speaker & \quad\quad\quad\quad\quad\quad\quad\quad\quad\quad\quad &
\textbf{English}  & C1-C2           \\
%\textbf{Spanish}  & A2-B1 & & \textbf{French}  & A2-B1
\\
\end{tabular}
\end{center}

\eatvspace

\begin{datetabular}{5em}
\dateentry{Dec 2023}{
\textbf{IELTS Academic Overall Band Score 8}

Obtained IELTS English certification, with CEFR level C1/C2.
}

\end{datetabular}


\section{Projects}


\begin{datetabular}{2em}


\dateentry{2023}{
\textbf{\href{https://github.com/luigi-pizza/ParallelPCA}{ParallelPCA}} 
\textit{Multicore Programming}

Designed and developed with a team of two other students, this project implements a distributed algorithm for Principal Components Analysis. ParallelPCA implements two different methodologies: multithreading and message passing between processes. 
}


\dateentry{2023}{
\textbf{\href{https://github.com/luigi-pizza/wasaPhoto}{WASAPhoto}} 
\textit{Full-Stack Web Application}

This project develops a social media platform that focuses on photo sharing. It implements a Go backend and Vue.js frontend.
}


\dateentry{2023}{
\textbf{\href{https://github.com/luigi-pizza/SelfDrivingDrone}{SelfDrivingDrone}} 
\textit{Computer Vision}

Designed and developed with a team of 2 other students, This project develops an autonomous obstacle navigation system for the DJI Tello drone. The drone is capable of autonomous flight through a predetermined set of obstacles. 
}



\dateentry{2022}{
\textbf{\href{https://github.com/luigi-pizza/BattleOfSexesSimulation}{Population simulator}} \textit{Java, JavaFX}

Designed and developed with a team of 2 other students, an evolution simulator in Java based on the 1976 book "The Selfish Gene" by Richard Dawkins, using multithreading to achieve a non-deterministic result. A GUI implemented with JavaFX enables the visualization of the simulation in real-time.
}


\end{datetabular}



% \section{Skills}
% \begin{datetabular}{15em}
% \dateentry{
% Programming Languages:

% Mathematics:

% Computer Architecture:

% AI: 

% Databases:

% Project Management Tools:}
% {C, C++, Java, Python

% Calculus, Multivariate Calculus, Linear Algebra, Probability, Statistics

% x86, RISC-V, Operative Systems, Networks

% Scikit-learn, PyTorch, Computer Vision, NLP

% SQL, MongoDB, Neo4J

% Git, Anaconda
% }
% \end{datetabular}


%\centerline{\textbf{Italian}  \quad Native Speaker \quad\quad\quad\quad\quad\quad\quad\quad
%\quad\quad\quad\quad\quad\quad\quad\quad\textbf{English}  \quad C1-C2}
%\\
%\quad\quad\quad\quad\quad\quad\quad\quad\quad\quad\textbf{Spanish}  \quad A2-B1 \quad\quad\quad\quad\quad\quad\quad\quad
%\quad\quad\quad\quad\quad\quad\quad\quad\quad\quad\textbf{French}  \quad A2-B1
%\\
%\section{Publications}
%\nocite{*} % Loads every entry from the attached .bib file
% Highlight takes three entries, given name, given name initials, and family name.
% If you have a middle initial, this call looks like:
% \highlightauthorname{Bob H.}{B. H.}{Smith}
%\highlightauthorname{Frodo}{F.}{Baggins} 
%\begin{datetabular}{2em}
%printbibyear has been defined to only print entries from a given citation year.
%\dateentry{3011}{\printbibyear{3001}}
%\dateentry{3001}{\printbibyear{3011}}
%\end{datetabular}

\end{document}
